\section{Controlled Intervention Details}\label{app:controlled-details}
The frozen controlled study uses \texttt{gpt-6-luna}, low reasoning, through the OpenAI Codex subscription interface. Each response is obtained from an ephemeral thread with no execution environment, available tools, workspace roots or inherited project instructions. Only synthetic public task inputs are transmitted. Original final answers, request/turn/usage records, model attestations and process-closure receipts accompany the results. Released transport projections omit account events and local identifiers, with hashes linking them to retained originals. There are 72 application admissions, 72 completed responses, zero application retries and 195,927 observed tokens. Collection takes 472.6 seconds. Provider-internal retries and underlying network-request counts are not exposed by this interface.

\paragraph{Receiver experiment}
The fixed seed 20261021 assigns 24 topic pairs to twelve training and twelve test topics; labels are brief and expanded private instructions. Each saved text traverses each arm once in a shuffled schedule, producing 144 trials and 288 observer/site views. TLS1.2 with ECDHE-RSA-AES128-GCM-SHA256 permits explicit parsing of application record types. The relays record encrypted byte streams, direction and local read timestamps; they cannot read the explanation. Four-byte application length headers are inside encryption. Padded frames have 1,024 application bytes and 24 bytes of GCM nonce/tag overhead, plus the five-byte TLS record header. Sixteen such records therefore cost 16,848 bytes. No padding-budget overflow or content-integrity failure occurs.

The attacker fits one length threshold and direction on training topics, including constant-prediction candidates with fixed tie-breaking. It is retrained independently for each protected arm/site. Held-out classifications are 24/24 at unpadded links and 12/24 at padded links. Upstream-only padding leaves terminal predictions unchanged; final-link padding changes terminal accuracy by $-50$ percentage points. Every paired test topic has the same difference, making the saved topic-bootstrap intervals degenerate. These intervals reflect this fixed workload and estimator; they do not establish a population-wide error bound. The definitive mechanism check is that the captured padded application-record feature is constant across all 48 outputs.

Median end-to-end replay time is 25.97ms without padding, 27.33ms with upstream-only padding and 28.52ms with both links padded. These descriptive local measurements include two TLS handshakes and frontend work, but exclude model generation and remote deployment latency. The wire multiplier is 808,704/34,606, approximately 23.37 per padded link across the 48 outputs. All 144 terminal content hashes match the source, and all owned socket threads join.

\paragraph{Execution experiment}
The 24 tasks split evenly between inventory and calendar domains. Twelve cover writes overlap the target resource; six use a disjoint protected resource and six use a disjoint spare resource. Three of the six overlapping inventory tasks have spare stock, while three have exactly the requested stock. All six overlapping calendar writes occupy the requested slot before the real booking. Cover writes have valid actor/confirmation parameters; failures therefore concern state semantics rather than malformed requests. All 24 model plans are valid and contain one write. No plan is dropped or repaired.

Across the five routes, 120 task/arm rows are retained as paired observations of 24 tasks. Live-write successes are 6/12 inventory and 3/12 calendar. Overlapping inventory succeeds in 3/6; overlapping calendar in 0/6. Each domain's disjoint cover succeeds on all six requested goals, but three protected-resource violations reduce full success to 3/6. This separates objective completion, hard constraints and exact state preservation. All four cover implementations retain their actual local tool call counts and elapsed time; local microseconds are not remote API costs.

The live-write full-success estimate is 37.5\% (Wilson 95\% interval 21.2--57.3\%), compared with 100\% (86.2--100\%) for the other routes. The prespecified paired statistic resamples whole tasks 10,000 times with seed 20261021 and gives a difference of $-62.5$ points, interval $[-79.2,-41.7]$. These intervals describe variation among constructed examples rather than a natural task population. The deterministic state-equality results apply to the implemented pure-read and whole-clone-isolation routes. Neither state equality nor utility implies indistinguishability to a privacy attacker.

\paragraph{Verification and retained outcomes}
Eighteen no-model tests cover protocol failures, missing usage, tool authorization and confirmation, malformed plans, state isolation, TLS parsing, content preservation and cleanup. The separate post-collection audit reads raw model events, ciphertext captures and before/after states and checks saved numerical results. It does not request new model responses or replace observations. Original empty/failed retail runs and the negative workflow-audit endpoint remain in Supplements S3/S7. New controlled results support the placement decisions for their explicit system and task constructions.
